\section{Analysis \& Discussion}

\subsection{Why Content Decays Faster Than Style}

Our central finding---that entities decay 5.8$\times$ faster than POS distributions as intensity rises---has a mechanistic explanation. POS distributions are \emph{statistical aggregates}: even aggressive paraphrasing that restructures every sentence tends to produce text with broadly similar grammatical composition; a noun phrase gets rephrased but typically remains a noun phrase.

Named entities, by contrast, are \emph{specific lexical items} that must be reproduced exactly to count as preserved. Minor reformulations (``New York City'' $\rightarrow$ ``NYC'' $\rightarrow$ ``the city'') register as entity loss under string matching, making NER a sharper probe than distributional metrics.

The POS tag delta heatmap reveals which specific tags drive the aggregate cosine shift. NOUN and PUNCT show the largest absolute changes under Dipper (High), consistent with text compression (shorter sentences, fewer proper nouns). Notably, PROPN (proper nouns) decreases across nearly all paraphrasers, aligning with the entity dropping we observe in the NER analysis.

Crucially, the multi-tier analysis confirms this is not a first-iteration artifact: the ordering BERTScore $>$ POS cosine $>$ NER Recall $>$ ROUGE-L $>$ BLEU holds at every tier, with no feature ``catching up'' to another. Erosion rates are intrinsic to feature type---aggregate distributional metrics (POS, BERTScore) are inherently more robust than point-specific ones (entities, exact n-grams).

\subsection{Paraphraser Fingerprints Are Multi-Dimensional}

The style-vs-content scatter (Figure~\ref{fig:fingerprints}) shows no single metric distinguishes all paraphrasers. Pegasus (Full) and Dipper (Default) share POS cosine ($\approx$0.963) but diverge on NER Precision (0.846 vs.\ 0.537): effective paraphraser identification (RQ3) requires \emph{multi-signal} classifiers.

\subsection{Implications for Forensic Applications}

Our findings have direct practical implications. First, plagiarism detection systems that rely solely on lexical overlap (e.g., n-gram matching) will fail after a single paraphrase iteration---BLEU drops 81\% at $T_1$ for ChatGPT. However, systems that incorporate POS distributions may remain effective longer, since grammatical structure persists even under aggressive paraphrasing (Dipper High retains 0.843 POS cosine at $T_3$).

Second, the layered erosion pattern suggests a defense-in-depth strategy for forensic analysis: check lexical overlap first (fast but fragile), then entity preservation (moderate sensitivity), then syntactic structure (most robust). Each layer provides evidence at a different sensitivity level, and the combination may enable reliable detection even after multiple paraphrase iterations.

Third, the finding that paraphrasers leave distinct multi-dimensional fingerprints is encouraging for RQ3. If different paraphrasers produce separable patterns in POS $\times$ NER $\times$ BERTScore space, classifiers trained on these features may be able to identify \emph{which} paraphraser was used, not just \emph{whether} paraphrasing occurred.

\subsection{The Ship of Theseus Revisited}

Returning to our framing metaphor: the grammatical skeleton (hull) of the ship is remarkably durable---even the most aggressive paraphraser retains 88\% cosine similarity in POS distribution at $T_1$, declining only to 84\% by $T_3$. The named entities (cargo) are jettisoned first and most dramatically. The lexical planks (individual words, measured by BLEU) are replaced almost entirely at the first iteration. Yet the ship's meaning (measured by BERTScore) persists throughout.

This layered erosion---lexical planks first, then entity cargo, with grammatical hull and semantic meaning largely intact---suggests that authorial identity is not a single property but a composite of features with different resilience profiles.

\subsection{Domain Effects}

Our domain-specific analysis reveals that decay patterns are not uniform across text types. Three factors explain the variation:

\textbf{Entity density drives content resilience.} XSum (BBC news) averages 18.7 entities per document and retains the most entities after paraphrasing (Recall 0.461 at $T_1$). Yelp (reviews) averages only 5.9 entities and shows the worst retention (0.362). Entity-dense texts provide more ``anchors'' that paraphrasers must explicitly handle, reducing the chance of inadvertent loss.

\textbf{Register formality drives style resilience.} CMV (argumentative essays) shows the highest POS cosine at $T_1$ (0.976), likely because argumentative writing uses conventional syntactic structures---subject-verb-object constructions, logical connectives, subordinate clauses---that paraphrasers reproduce naturally. Yelp reviews, with more colloquial and variable syntax, offer paraphrasers more room to restructure.

\textbf{Attribution robustness varies by domain.} Authorial identity persists longest in CMV (F1 = 0.706 at $T_1$) and news (XSum: 0.632), but collapses toward chance in explanatory text (ELI5: 0.542). This has practical implications: forensic authorship tools may need domain-specific calibration, as the ``point of no return'' is domain-dependent.

\subsection{Source Author Effects}

The finding that human-authored texts show slightly less POS drift than LLM-generated texts, while counterintuitive, is consistent with the observation that six distinct LLM architectures produce syntactically heterogeneous $T_0$ texts. A single paraphraser normalizes these diverse inputs toward its own syntactic template, producing measurable POS shifts. Human writing, being more syntactically conventional, offers less room for restructuring. This has a practical implication: paraphrasing detectors may perform differently on human-authored vs.\ LLM-authored source texts, and evaluation benchmarks should account for this asymmetry.
